\section{Discussion}
\label{sec:discussion}

\para{Why did the framing fail?} The hypothesis imagines full-rank matrices that collapse. With standard initialization, scaled dot-product attention begins near-uniform (row entropy about 0.99) because queries and keys are small, so $\mA$ is close to a rank-one matrix and learning builds structure. The $\times3$ arm is the key control. When attention starts peaked (row entropy 0.48, $\reff/n$ about 0.32 on \agnews), rank falls, and falls more with depth. Last layers then end close together from either side, at .06 for default and .09 for $\times3$ on \agnews, though first layers keep their starting level (.07 and .32). The start state therefore persists early in the network and matters less late in it. \bert starts peaked as well (row entropy about 0.50), and its last layers lose rank during fine-tuning, which agrees with this reading.

\para{Practical implications.} We advise against a fixed $0.3\log n$ cut-off for pruning. Descriptions of attention training as rank collapse should state the initial state, since it sets the direction. On vision models, entropy is a useful ranking signal, but this does not carry over to text. We did not expect this split. One possible reason is that a fixed image grid gives heads stable spatial roles while text attention varies by input, but we did not test this.

\para{Alternative explanations.} {\it Small models:} \bert is far larger and pretrained, and it is also unimodal and mostly low-rank (mean dip $p=.99$). {\it Pre-LN:} post-LN matched pre-LN almost exactly. {\it \textsc{cls} pooling:} this was constant across arms, so we cannot rule out an effect. {\it Averaging across inputs:} per-head means could hide bimodality at the level of single inputs, and this is a real limitation.

\para{Limitations.} We study small encoders and short fine-tunes on three classification datasets. We did not study ImageNet-scale vision transformers, decoders, or large language models, where bifurcation might hold, although nothing here suggests it. The definitions of $\reff$ and $S$ were our choice, as the hypothesis gave none, and five rank definitions cover only H1. The $\times3$ arm, rise-time ratios, $\epsilon=2$ as reporting scale, and paired baseline tests were not pre-registered, and we tested only one gain value, so these analyses are exploratory.

For H3, \sst evaluation has 436 sentences, \bert has 3 seeds with very wide intervals, and we ran fine-tuning only on $L=6$ default arms. Zeroing a head without rescaling is one of several removal methods.

For H4, noise channels define a quantity only relative to $\epsilon$. Layer-level and per-head values differ by about a factor of eight at a fixed $\epsilon\ge1$, and the claim does not say which it means. \agnews's intrinsic $H(Y|X)$ is assumed near zero and is really some small unknown value. Injected noise applies to training labels, while clean test inputs supply $X$. On \order, training failed at small $n$ (test accuracy near 50\% at $n=500$ and for noisy runs at $n=2{,}000$), and 3-seed means are noisy. Our reading of $n$ as training-set size is an assumption. A derivation with stronger assumptions, such as stochastic attention or a different left side such as $I(\mA;Y)$, could survive, but it would not be this claim.

Time to minimum uses a 5\% tolerance on 41 checkpoints, so differences below about 1\% of training early on, or 4\% later, cannot be resolved. Shared-GPU restarts should not affect results, and the exact rerun supports this, but it covered one configuration on one GPU type.
